\documentclass{article}
\usepackage{hyperref}
\usepackage{Style}

\nocite{*} % Comentar si quiero citar
%\addbibresource{bibliografia.bib} % Quitar el comentado si quiero usar bibliografia

\begin{document}

\begin{minipage}{2.5cm}
    \includegraphics[width=2cm]{imagen_puc.jpg}
\end{minipage}
\begin{minipage}{12cm}
    {\sc Pontificia Universidad Católica de Chile\\
    Facultad de Matemáticas\\
    Departamento de Matemática\\
    Profesor: Manuel Cabezas -- Ayudante: Benjamín Mateluna}
\end{minipage}
\vspace{2ex}

{\centerline{\bf Introducción a la Geometría - MAT1304}
\centerline{\bf Ayudantía 5}}
\centerline{\bf 23 de marzo de 2026}

\vspace{1cm}
\noindent\textbf{Problema 1.} Dado $\triangle ABC$, consideremos $Z$ un punto en $\overline{AB}$.
Se traza el segmento $\overline{CZ}$, trazamos la recta paralela a $\overline{CZ}$ que intersecta
a $\overleftrightarrow{BC}$ en $X$. Del mismo modo, se traza la recta paralela a $\overline{CZ}$ 
que pasa por $B$ y que intersecta a $\overleftrightarrow{AC}$ en $Y$. Muestre que
\begin{equation*}
    \frac{1}{AX}+\frac{1}{BY}=\frac{1}{CZ}
\end{equation*}

\vspace{5mm}
\noindent\textbf{Problema 2.} Sea $\triangle ABC$ un triángulo cualquiera. Consideremos los puntos 
$D,E,F$ en $\overline{BC},\overline{AC},\overline{AB}$ respectivamente, tales que 
$\overline{AD},\overline{BE}$ y $\overline{CF}$ concurren dentro del triángulo. Muestre que 
si $\overline{FE}$ y $\overline{BC}$ son paralelas, entonces $\overline{BD}=\overline{DC}$.

\vspace{5mm}
\noindent\textbf{Problema 3.} Sea $\triangle ABC$ un triángulo cualquiera. Sean $D,E$ en 
$\overline{AB}$ y $\overline{BC}$ respectivamente tales que los segmentos $\overline{AE}$ y 
$\overline{CD}$ son alturas. Muestre que $\triangle ABE\sim\triangle CBD$.

\vspace{5mm}
\noindent\textbf{Problema 4.} Sean ABCF un cuadrado y $\triangle CDE$ un triángulo equilátero,
ambos de lado 2 y tales que $B,C$ y $D$ son colineales. Sean $H\in\overline{AD}\cap\overline{EC}$ 
y $G\in\overline{AD}\cap\overline{FC}$. Calcular el área del triángulo $\triangle GCH$.

\vspace{1cm}
\noindent\textbf{\underline{Ejercicios Propuestos:}}

\vspace{5mm}
\noindent\textbf{Extra 1.} Dados $\triangle ABC$ y $\triangle DEF$ semejantes con razón $k$, 
pruebe que
\begin{equation*}
    \frac{\abs{\triangle ABC}}{\abs{\triangle DEF}}=k^{2}
\end{equation*}

\vspace{5mm}
\noindent\textbf{Extra 2.} Considere la figura de abajo. Demuestre que 
$\overline{AC}>\overline{DE}$.

\begin{center}
    \includegraphics[scale=0.3]{E2.png}
\end{center}

%\printbibliography % Quitar el comentado si quiero usar bibliografia

\end{document}